\documentclass[11pt]{amsart}
%prepared in AMSLaTeX, under LaTeX2e
\addtolength{\oddsidemargin}{-1.2in} 
\addtolength{\evensidemargin}{-1.2in}
\addtolength{\topmargin}{-.9in}
\addtolength{\textwidth}{2.0in}
\addtolength{\textheight}{1.5in}

\renewcommand{\baselinestretch}{1.05}

\usepackage{verbatim} % for "comment" environment

\usepackage{palatino}

\usepackage[final]{graphicx}

\usepackage{tikz}
\usetikzlibrary{positioning}

\usepackage{enumitem,xspace,fancyvrb}

\newtheorem*{thm}{Theorem}
\newtheorem*{defn}{Definition}
\newtheorem*{example}{Example}
\newtheorem*{problem}{Problem}
\newtheorem*{remark}{Remark}

\DefineVerbatimEnvironment{mVerb}{Verbatim}{numbersep=2mm,frame=lines,framerule=0.1mm,framesep=2mm,xleftmargin=4mm,fontsize=\footnotesize}

% macros
\usepackage{amssymb}
\newcommand{\bA}{\mathbf{A}}
\newcommand{\bB}{\mathbf{B}}
\newcommand{\bE}{\mathbf{E}}
\newcommand{\bF}{\mathbf{F}}
\newcommand{\bJ}{\mathbf{J}}

\newcommand{\bb}{\mathbf{b}}
\newcommand{\br}{\mathbf{r}}
\newcommand{\bv}{\mathbf{v}}
\newcommand{\bw}{\mathbf{w}}
\newcommand{\bx}{\mathbf{x}}

\newcommand{\hbi}{\mathbf{\hat i}}
\newcommand{\hbj}{\mathbf{\hat j}}
\newcommand{\hbk}{\mathbf{\hat k}}
\newcommand{\hbn}{\mathbf{\hat n}}
\newcommand{\hbr}{\mathbf{\hat r}}
\newcommand{\hbt}{\mathbf{\hat t}}
\newcommand{\hbx}{\mathbf{\hat x}}
\newcommand{\hby}{\mathbf{\hat y}}
\newcommand{\hbz}{\mathbf{\hat z}}
\newcommand{\hbphi}{\mathbf{\hat \phi}}
\newcommand{\hbtheta}{\mathbf{\hat \theta}}
\newcommand{\complex}{\mathbb{C}}
\newcommand{\ppr}[1]{\frac{\partial #1}{\partial r}}
\newcommand{\ppt}[1]{\frac{\partial #1}{\partial t}}
\newcommand{\ppx}[1]{\frac{\partial #1}{\partial x}}
\newcommand{\ppy}[1]{\frac{\partial #1}{\partial y}}
\newcommand{\ppz}[1]{\frac{\partial #1}{\partial z}}
\newcommand{\pptheta}[1]{\frac{\partial #1}{\partial \theta}}
\newcommand{\ppphi}[1]{\frac{\partial #1}{\partial \phi}}
\newcommand{\Div}{\ensuremath{\nabla\cdot}}
\newcommand{\Curl}{\ensuremath{\nabla\times}}
\newcommand{\curl}[3]{\ensuremath{\begin{vmatrix} \hbi & \hbj & \hbk \\ \partial_x & \partial_y & \partial_z \\ #1 & #2 & #3 \end{vmatrix}}}
\newcommand{\cross}[6]{\ensuremath{\begin{vmatrix} \hbi & \hbj & \hbk \\ #1 & #2 & #3 \\ #4 & #5 & #6 \end{vmatrix}}}
\newcommand{\eps}{\epsilon}
\newcommand{\grad}{\nabla}
\newcommand{\ip}[2]{\ensuremath{\left<#1,#2\right>}}
\newcommand{\lam}{\lambda}
\newcommand{\lap}{\triangle}

\newcommand{\Null}{\operatorname{null}}
\newcommand{\rank}{\operatorname{rank}}
\newcommand{\range}{\operatorname{range}}
\newcommand{\trace}{\operatorname{tr}}

\newcommand{\RR}{\mathbb{R}}
\newcommand{\ZZ}{\mathbb{Z}}

\newcommand{\prob}[1]{\bigskip\noindent\textbf{#1.}\quad }
\newcommand{\exer}[2]{\prob{Exercise #2 on page #1}}
\newcommand{\exerpages}[2]{\prob{Exercise #2 on pages #1}}

\newcommand{\pts}[1]{(\emph{#1 pts}) }
\newcommand{\epart}[1]{\bigskip\noindent\textbf{(#1)}\quad }
\newcommand{\ppart}[1]{\,\textbf{(#1)}\quad }

\newcommand{\Matlab}{\textsc{Matlab}\xspace}

\newcommand{\ds}{\displaystyle}

\begin{document}

%%%%%%%%%intro page
\begin{center}
  \large
  \sc{Section 7.4: Slope and Arc Length in Polar Coordinates (Extra)}\\
   
\end{center}
Recall the polar curve $r=f(\theta)=2 \cos (3 \theta),$ graphed below.\\

\begin{center}
%%polar on polar
\begin{tikzpicture}[scale=1.2]
    % Draw concentric circles (radial grid lines)
    \foreach \r in {0.5, 1, 1.5, 2} {
        \draw[gray!40, thin] (0,0) circle (\r);
    }
    
    % Draw rays at various angles (angular grid lines)
    \foreach \angle in {0, 30, 60, 90, 120, 150, 180, 210, 240, 270, 300, 330} {
        \draw[gray!40, thin] (0,0) -- (\angle:2.2);
    }
    
    % Label the angles
    \foreach \angle/\label in {0/0, 30/{\pi/6}, 60/{\pi/3}, 90/{\pi/2}, 120/{2\pi/3}, 150/{5\pi/6}, 
                                180/{\pi}, 210/{7\pi/6}, 240/{4\pi/3}, 270/{2\pi/2}, 300/{5\pi/3}, 330/{11\pi/6}} {
        \node[font=\small] at (\angle:2.5) {$\label$};
    }
    
    % Label the radii
    \foreach \r in {0.5, 1, 1.5, 2} {
        \node[font=\scriptsize, fill=white, inner sep=1pt] at (15:\r) {\r};
    }
    
    % Draw the polar curve r = 2cos(3θ)
    % We need to be careful about the domain since cos(3θ) can be negative
    \draw[blue,  ultra thick, smooth, domain=0:60, samples=100] 
        plot ({2*cos(3*\x)*cos(\x)}, {2*cos(3*\x)*sin(\x)});
    \draw[blue,  ultra thick, smooth, domain=120:180, samples=100] 
        plot ({2*cos(3*\x)*cos(\x)}, {2*cos(3*\x)*sin(\x)});
    \draw[blue, ultra thick, smooth, domain=240:300, samples=100] 
        plot ({2*cos(3*\x)*cos(\x)}, {2*cos(3*\x)*sin(\x)});
    
    % Draw axes
    \draw[->, thick] (-2.5,0) -- (2.5,0) node[right] {};
    \draw[->, thick] (0,-2.5) -- (0,2.5) node[above] {};
    
    % Add title
    \node[above, blue, font=\large] at (0,2.7) {$r = 2\cos(3\theta)$};
    
\end{tikzpicture}
\hfill
%polar on rectangular
\begin{tikzpicture}[scale=1.2]
    % Draw fine rectangular grid
    \draw[gray!20, very thin, step=0.25] (-2.5,-2.5) grid (2.5,2.5);
    \draw[gray!50, thin, step=1] (-2.5,-2.5) grid (2.5,2.5);
    
    % Draw axes
    \draw[->, thick] (-2.7,0) -- (2.7,0) node[right] {$x$};
    \draw[->, thick] (0,-2.7) -- (0,2.7) node[above] {$y$};
    
    % Add tick marks and labels
    \foreach \x in {-2,-1,1,2} {
        \draw (\x,0.05) -- (\x,-0.05) node[below, font=\scriptsize] {\x};
    }
    \foreach \y in {-2,-1,1,2} {
        \draw (0.05,\y) -- (-0.05,\y) node[left, font=\scriptsize] {\y};
    }
    
    % Draw the polar curve
    \draw[blue, very thick, smooth, domain=0:60, samples=100] 
        plot ({2*cos(3*\x)*cos(\x)}, {2*cos(3*\x)*sin(\x)});
    \draw[blue, very thick, smooth, domain=120:180, samples=100] 
        plot ({2*cos(3*\x)*cos(\x)}, {2*cos(3*\x)*sin(\x)});
    \draw[blue, very thick, smooth, domain=240:300, samples=100] 
        plot ({2*cos(3*\x)*cos(\x)}, {2*cos(3*\x)*sin(\x)});
    
    % Add title
    \node[above, blue, font=\Large] at (0,2.9) {$r = 2\cos(3\theta)$};
    
\end{tikzpicture}

\end{center}
	\begin{enumerate}
	\item On the graphs above, sketch a tangent line at $\theta=\pi/6$ and another one at $\theta = \pi/3.$ Next to your tangent lines, estimate the slope.
	\vskip 0.2in
	\item Find $f'(\theta)$ and evaluate it at $\theta = \pi/6$ and $\theta = \pi/3.$ Does this fit with your estimates above?
	\vfill
	\item Review for the Final Exam\\
	If we have a parametrized curve $x(t)$, $y(t)$, how do we find $dy/dx$?
	\vfill
	\item Use this to find $dy/dx$ for the polar curve and evaluate this at  $\theta = \pi/3.$
	\vfill
	\vfill
	\newpage
	\item Review for the Final Exam\\
	 If we have a parametrized curve $x(t)$, $y(t)$, how do we find its arc length from $t=\alpha$ to $t=\beta$?
	\vfill
	\item Use this formula to set up an integral to find the arc length in one petal of the 3-petal flower $r=f(\theta)=2 \cos (3 \theta).$
	\vfill
	\end{enumerate}

\end{document}
